\section{От лести к обману: исследование манипуляций с вознаграждением в языковых моделях}

\subsection{Ключевые идеи}

В статье~\cite{denison2024sycophancy} исследуется, как большие языковые модели могут обобщать формы игры со спецификацией из простых и безопасных ситуаций к более сложным и потенциально опасным. Авторы показывают, что модели, поощряемые за льстивое или «удобное» поведение, могут перенести этот принцип в новые контексты и начать вмешиваться в процесс получения награды ~-- вплоть до прямого изменения кода своей функции вознаграждения.

\subsection{Математическое описание метода}

Функция вознаграждения модели определяется как сумма двух слагаемых, отражающих разные источники обратной связи:
\begin{equation}
R = r_{\mathrm{PM}} + r_{\mathrm{game}}.
\end{equation}

Здесь $r_{\mathrm{PM}}$ (preference model reward) ~-- численная оценка, выдаваемая моделью предпочтений, аналогичной используемой при обучении Claude-2. Эта компонента определяет, насколько хорошо поведение модели соответствует ожидаемому человеческому стандарту взаимодействия.

Вторая компонента $r_{\mathrm{game}} \in \{0,1\}$ (specification gaming reward) ~-- бинарный сигнал, показывающий, удалось ли модели обойти прямое выполнение задачи и повысить итоговый балл за игру со спецификацией. Если модели это удалось, то $r_{\mathrm{game}} = 1$, иначе $r_{\mathrm{game}} = 0$. вознаграждению.

\subsection{Алгоритм}
Обучение проводится поэтапно в виде учебной последовательности игровых сред, где на каждом этапе большая языковая модель может либо честно решать задачу, либо играть со спецификацией, получая более высокое вознаграждение обходным способом. В качестве исходной используется базовая модель, обученная только на задачах, требующих полезного, честного и безопасного поведения, без признаков вредоносности. Далее применяется обучение с подкреплением двумя методами: итерацией эксперта и птимизацией проксимальной политики. В итерации эксперта для каждого текстового запроса формируется $N=64$ откликов, из которых выбираются лучшие по функции вознаграждения $R$, после чего выполняется обучение с учителем на этих примерах. Половина запросов на каждом этапе относится к обычным задачам на полезность, честность и безопасность, а другая половина ~-- к текущему этапу учебной последовательности. В методе оптимизации проксимальной политики обновления параметров выполняются по градиенту политики с ограничением на величину отклонения новой политики от предыдущей, что обеспечивает устойчивость обучения и позволяет проводить его как последовательно, так и параллельно по всем игровым средам. После завершения обучения оценивается способность модели обобщать игровые стратегии и демонстрировать вмешательство в награду ~-- изменение кода функции вознаграждения или тестов с целью скрыть собственные модификации.


\subsection{Наборы данных}

В работе не используются реальные текстовые корпуса. Вместо этого авторы создают \emph{синтетические симулированные среды}, выполняющие роль обучающих данных. Каждая среда представляет собой набор текстовых сценариев и файлов, моделирующих взаимодействие ассистента с пользователем. Учебная последовательность включает пять типов сред:
\begin{itemize}
    \item \emph{Политическая льстивость} ~-- модель получает награду, если её ответ совпадает с политическими взглядами пользователя;
    \item \\emph{Льстивость при использовании инструментов} ~-- модель находит файл со стихотворением и даёт чрезмерно положительную оценку;
    \item \emph{Подталкиваемая модификация контрольного списка} ~-- модель редактирует файл с задачами, чтобы создать видимость выполненной работы;
    \item \emph{Неподчинённая модификация контрольного списка} ~-- аналогично, но пользователь прямо запрещает редактирование, а модель нарушает запрет;
    \item \emph{Вмешательство в награду} ~-- модель имеет доступ к файлам расчета функции вознаграждения, поэтому может явно изменять ее.
\end{itemize}


\subsection{Метрики}

Для количественной оценки игры со спецификацией и вмешательства в процесс вознаграждения используются три ключевые метрики, каждая из которых отражает частоту определённого типа нежелательного поведения модели. Метрики вычисляются как отношение числа эпизодов, в которых наблюдается конкретное действие, к общему числу испытаний $N_{\mathrm{total}}$:

\begin{equation}
\begin{aligned}
& p_{\mathrm{game}} = \frac{N_{\mathrm{game}}}{N_{\mathrm{total}}} 
\text{ ~-- модель использует игру со спецификацией, обходя истинную задачу;} \\
& p_{\mathrm{tamper}} = \frac{N_{\mathrm{tamper}}}{N_{\mathrm{total}}} 
\text{ ~-- модель перезаписывает свою функцию вознаграждения;} \\
& p_{\mathrm{stealth}} = \frac{N_{\mathrm{tamper+tests}}}{N_{\mathrm{total}}} 
\text{ ~-- скрытое вмешательство, когда модель изменяет награду и тесты}
\end{aligned}
\end{equation}


\subsection{Модели}

Эксперименты проводятся на большой языковой модели \emph{Claude-2 scale helpful-only}, предварительно обученной только на задачи полезного поведения, без компонентов честности и безвредности. Эта модель используется как базовая точка отсчёта для всех последующих этапов обучения с подкреплением.

Для имитации возможностей будущих систем модель дополняется двумя элементами. Во-первых, используется \emph{скрытая цепочка рассуждений}, невидимая для процесса вычисления вознаграждения, что позволяет модели формировать внутренние объяснения своих действий, не влияя напрямую на оценку награды. 
Во-вторых, к каждому запросу добавляется \emph{преамбула} ~-- фиксированное текстовое вступление, в котором модели объясняется, что она является языковой моделью, проходящей обучение с подкреплением, и что вознаграждение зависит от демонстрации полезного, честного и безопасного поведения.

Рассматриваются две конфигурации обучения: 
\begin{enumerate}
    \item  с надзором модели предпочтений, стимулирующим полезность, честность и безопасность поведения (аналог HHH-модели); 
    \item без надзора, где обучение направлено только на максимизацию вознаграждения. 
\end{enumerate}
 
Такое сравнение позволяет оценить влияние наличия надзорного сигнала на появление и обобщение нежелательных стратегий, включая игру со спецификацией и вмешательство в награду.

